\documentclass[preprint,12pt]{elsarticle}
\myfooterfont{\scriptsize\itshape}
\makeatletter
\patchcmd{\ps@pprintTitle}
  {\hbox to \textwidth}
  {\hbox to \dimexpr\textwidth-3pt\relax}
  {}{}
\makeatother

\usepackage{amsmath,amssymb}
\usepackage{booktabs}
\usepackage{graphicx}
\usepackage[hidelinks]{hyperref}
\usepackage{siunitx}
\usepackage{xcolor}
\usepackage{lineno}
\setlength{\emergencystretch}{2em}
\makeatletter
\pdfstringdefDisableCommands{%
  \def\corref#1{}%
  \def\cormark[#1]{}%
  \def\cnotenum{}%
  \def\@corref{}%
}
\setlength{\@fptop}{0pt}
\setlength{\@fpbot}{0pt plus 1fil}
\makeatother

\journal{International Journal of Heat and Mass Transfer}
\newcommand{\OpenFOAM}{OpenFOAM}

\begin{document}
\begin{frontmatter}

\title{Supplementary material for ``Transient conjugate heat transfer in internally heated ceramic pebble beds: pore-resolved reference fields and graph--Transformer prediction''}

\author[a,b]{Jian Wang\corref{cor1}}
\ead{wjfttt@mail.ustc.edu.cn}
\author[a,b]{Wei Wen}
\author[a]{Gang Shen}
\author[b]{Mingzhun Lei}
\author[b]{Yuntao Song}
\cortext[cor1]{Corresponding author.}

\address[a]{Anhui University of Science and Technology, Huainan 232001, China}
\address[b]{Institute of Plasma Physics, Hefei Institutes of Physical Science, Chinese Academy of Sciences, Hefei 230031, China}

\begin{abstract}
This supplementary material records the predefined data separation, the
legacy unbounded-output diagnostic and the numerical sensitivity results that
support the main paper. Complete inputs, programs and processed data are
provided in the reproduction package rather than repeated as multi-page
tables.
\end{abstract}

\end{frontmatter}
\linenumbers

\section{Numerical details and data separation}

The \OpenFOAM{} fields use the literature-based geometry, operating conditions
and thermophysical relations summarized in the main paper. The complete
machine-readable parameter table, equation-to-input map, numerical settings
and result-generating scripts are provided in the associated reproduction package
rather than repeated here as multi-page tables.

Steady field labels denote nonlinear iterations. Each accepted condition
reaches iteration 200 and is checked for finite outlet temperature, pressure
drop, maximum solid temperature, wall heat, mass difference and energy
difference. The fixed-flow trajectories use identical staged output times.
Grid- and time-step convergence indices are reported only for monotonic
three-level sequences; raw values are retained otherwise.

\begin{table}[!htbp]
\centering
\scriptsize
\setlength{\tabcolsep}{2.5pt}
\caption{Pre-defined physical data splits. A steady condition is one complete three-dimensional field at a fixed inlet velocity, inlet temperature and pebble heat source. A transient sample is one complete staged-time trajectory; no time point from a trajectory is assigned to another role. The endpoint-pair split also keeps the forward and reverse trajectories joining the same two steady conditions together. Normalization and checkpoint selection use training and validation data only. Exact condition and trajectory identifiers are stored in the cited machine-readable split files.}
\label{tab:physical_data_splits}
\begin{tabular}{>{\raggedright\arraybackslash}p{2.6cm}>{\raggedright\arraybackslash}p{3.2cm}>{\raggedright\arraybackslash}p{3.2cm}>{\raggedright\arraybackslash}p{3.2cm}}
\toprule
Physical test & Training & Validation & Independent test \\
\midrule
\multicolumn{4}{l}{\textit{Steady three-dimensional fields}} \\
Interleaved combinations & all sampled levels (36) & all sampled levels (12) & all sampled levels (12) \\
Temperature extrapolation & $T_{\rm in}=300, 500$ K (30) & $T_{\rm in}=700$ K (15) & $T_{\rm in}=900$ K (15) \\
Velocity extrapolation & $u_{\rm in}=0.05, 0.1, 0.15$ m s$^{-1}$ (36) & $u_{\rm in}=0.2$ m s$^{-1}$ (12) & $u_{\rm in}=0.25$ m s$^{-1}$ (12) \\
Heat-source interpolation & $q'''=4.85, 8.85$ MW m$^{-3}$ (32); disjoint combinations & $q'''=4.85, 8.85$ MW m$^{-3}$ (8); disjoint combinations & $q'''=6.85$ MW m$^{-3}$ (20) \\
Heat-source extrapolation & $q'''=4.85$ MW m$^{-3}$ (20) & $q'''=6.85$ MW m$^{-3}$ (20) & $q'''=8.85$ MW m$^{-3}$ (20) \\
\midrule
\multicolumn{4}{l}{\textit{Complete thermal-step trajectories}} \\
Downward-step prediction & six paired curves (6) & three upward curves (3) & three downward curves (3) \\
Upward-step prediction & six paired curves (6) & three downward curves (3) & three upward curves (3) \\
Endpoint-pair prediction & three endpoint pairs (6) & one endpoint pair (2) & two unseen endpoint pairs (4) \\
\bottomrule
\end{tabular}
\end{table}


No cell from an independent steady condition and no time point from an
independent trajectory is used for fitting, normalization or model selection.
Both directions connecting the same endpoint pair remain in one role. Model
architecture and normalization are frozen before the high-velocity and
independent-packing tests.

\subsection{Legacy unbounded-output diagnostic}

An early graph--Transformer could leave the thermophysical temperature range.
Those checkpoints are retained only to document why all final neural models
use the common differentiable bounded output map; they are excluded from the
accuracy comparison.

\IfFileExists{generated_unbounded_output_diagnostic.tex}{%
\begin{table}[htbp]
\centering
\scriptsize
\setlength{\tabcolsep}{3pt}
\caption{Diagnostic results from superseded graph--Transformer checkpoints with unbounded temperature outputs. These checkpoints are excluded from the final model comparison. The pressure channel remained positive, whereas the learned temperature channels left the specified thermophysical intervals; the physics-constrained checkpoint produced four non-positive fluid-temperature values.}
\label{tab:unbounded_output_diagnostic}
\begin{tabular}{@{}lrrrrr@{}}
\toprule
Checkpoint & Epochs & Min. $T_f$ (K) & Min. $T_s$ (K) & Min. $p_{\mathrm{abs}}$ (Pa) & $T_f\leq0$ \\
\midrule
Data-only & 500 & 94.524 & 69.493 & 119999.4 & 0 \\
Physics-constrained & 42 & -18.023 & 318.144 & 119999.4 & 4 \\
\bottomrule
\end{tabular}
\end{table}
}{}

\section{Additional checks and data access}

The independent-packing comparison and the PREMUX, TESOMEX, HELOKA and
fixed-bed comparisons are shown once in the main paper and are not duplicated
here. Their pointwise processed data, uncertainty fields and generating
programs are included in the reproduction package. The external measurements
check aggregate thermal and hydraulic scales; they are not cellwise labels or
a three-dimensional validation of the present local packing.

The repository contains the complete physical-parameter table, implementation
map, model settings, fixed split files, processed results and figure scripts.
It also contains the regional DMDc field comparison used to check the classical
baseline without repeating a second temperature-cloud figure in the paper.
Raw \OpenFOAM{} fields are archived separately because of their size. The paper
and supplementary material are regenerated with
\texttt{run\_hccb\_p418\_manuscript\_refresh.sh}.

\section{Numerical sensitivity}

\IfFileExists{generated_mesh_sensitivity.tex}{%
\begingroup
\renewcommand{\cite}[1]{Celik et al.\ (2008; full reference in the main paper)}
% Generated by code/build_hccb_p418_mesh_sensitivity_table.py
\begin{table*}[htbp]
\centering
\caption{Three-mesh sensitivity for the representative conjugate-heat-transfer condition. The equivalent cell size is $[(V_f+V_s)/(N_f+N_s)]^{1/3}$. Grid Convergence Index (GCI) values use the actual unequal refinement ratios and a safety factor of 1.25 \cite{celik2008gci}; no percentage GCI is reported for an oscillatory or unresolved sequence, or when the response crosses zero.}
\label{tab:mesh_sensitivity}
\small
\setlength{\tabcolsep}{4pt}
\resizebox{\textwidth}{!}{%
\begin{tabular}{@{}lrrrrrrr@{}}
\toprule
Grid & $N$ & $h_{\rm eq}$ ($\mu$m) & $\phi$ & $\Delta p$ (Pa) & $\Delta T_{\rm out}$ (K) & $\Delta T_{s,\max}$ (K) & $Q_{\rm wall}/Q_{\rm gen}$ \\
\midrule
Coarse & 361151 & 54.38 & 0.38591 & 23.3 & -26.52 & -0.1446 & -0.2733 \\
Medium & 947924 & 39.42 & 0.38661 & 25.64 & -26.2 & 0.05209 & -0.2942 \\
Fine & 1870064 & 31.43 & 0.38680 & 27 & -26.05 & 0.1608 & -0.3079 \\
\bottomrule
\end{tabular}
}
\par\vspace{0.5ex}
\resizebox{\textwidth}{!}{%
\begin{tabular}{@{}lrrr@{}}
\toprule
Quantity & Refinement behavior & Observed order & Fine-grid GCI \\
\midrule
Pressure drop & monotonic & 0.692 & 37.153\% \\
Outlet temperature rise & monotonic & 1.387 & 1.962\% \\
Maximum solid-temperature change & near-zero sign change & -- & -- \\
Signed wall heat / generated power & monotonic & 0.274 & 86.657\% \\
\bottomrule
\end{tabular}
}
\end{table*}

\endgroup
}{}

\IfFileExists{generated_timestep_sensitivity.tex}{%
\begin{table*}[htbp]
\centering
\footnotesize
\setlength{\tabcolsep}{5pt}
\caption{Temporal-resolution comparison for the representative heat-source step. The computed trajectories use the finest staged schedule: \num{1e-05}\,s over \num{0}--\num{0.1}\,s, \num{0.0005}\,s over \num{0.1}--\num{1}\,s, \num{0.01}\,s over \num{1}--\num{25}\,s. The second column reports the largest next-coarser--finest curve difference normalized by the finest-curve response span. GCI values are reported only for monotonically converging three-resolution triplets.}
\label{tab:thermal_timestep}
\begin{tabular}{lccc}
\toprule
Quantity & \shortstack{Finest-pair curve\\difference} & \shortstack{Endpoint\\GCI} & \shortstack{Curve-maximum\\GCI} \\
\midrule
Outlet temperature & 0.0385\% & 3.02e-05\% & 3.02e-05\% \\
Cooling-wall power & 0.0276\% & 0.000152\% & 0\% \\
Maximum solid temperature & 0.0425\% & 1.49e-05\% & 1.49e-05\% \\
Volume-averaged fluid temperature & 0.0359\% & 2.02e-05\% & 2.02e-05\% \\
Volume-averaged solid temperature & 0.0356\% & 2.03e-05\% & 2.03e-05\% \\
\bottomrule
\end{tabular}
\end{table*}
}{}

\IfFileExists{generated_final_uncertainty.tex}{%
\input{generated_final_uncertainty}}{}


\end{document}
